% This document was generated by an LLM.

\documentclass{essay}

\newcommand{\ddx}[1]{\frac{\dd #1}{\dd x}}
\newcommand{\pd}[2]{\frac{\partial #1}{\partial #2}}

\title{Differentials, Demystified}
\subtitle{Why you \emph{can} multiply by $\dd x$ in a separable ODE,
even though $\dfrac{\dd y}{\dd x}$ is ``not a fraction''}
\author{}
\date{}

\begin{document}

\maketitle

\textbf{The paradox, stated plainly.}
You have been told two things that seem to contradict each other:
\begin{enumerate}[label=\textbf{(\arabic*)}]
  \item $\dfrac{\dd y}{\dd x}$ is \emph{not} a fraction. It is a
    single symbol for a limit, so you are not allowed to split it
    into a numerator $\dd y$ and a denominator $\dd x$.
  \item Yet to solve a separable equation you write $\dd y =
    f(x)\,\dd x$, integrate both sides, and it always gives the right answer.
\end{enumerate}
Both statements are correct. The resolution is that step~(2) is a
\textbf{shorthand}: every ``illegal'' move is silently invoking a
genuine theorem that happens to make the notation behave like a
fraction. This document expands the shorthand back into the
theorems it stands for.

\section{First, why the warning is literally true}

By definition,
\[
  \ddx{y} \;=\; \lim_{\Delta x \to 0}\frac{\Delta y}{\Delta x},
\]
where $\Delta y = y(x+\Delta x) - y(x)$. The quotient $\Delta y /
\Delta x$ \emph{is} an honest fraction of two real numbers. But the
derivative is the \emph{limit} of that fraction, and a limit is a
single number. There is no leftover ``$\dd y$'' or ``$\dd x$''
sitting around once the limit has been taken --- they were consumed
in the limiting process. So at the level of introductory calculus,
$\dd y/\dd x$ is one indivisible symbol, and the warning is right.

The interesting question is therefore not ``is it a fraction?'' (it
is not) but ``\emph{why does treating it like one keep working?}''
The answer is a single theorem.

\section{The one theorem behind all of it}

\begin{theorem}[Substitution rule, i.e.\ the chain rule run backwards]
  If $F' = f$ and $u$ is a differentiable function of $x$, then
  \[
    \int f\bigl(u(x)\bigr)\,u'(x)\,\dd x \;=\; F\bigl(u(x)\bigr)+C
    \;=\; \int f(u)\,\dd u .
  \]
\end{theorem}
The famous mnemonic ``\,$\dd u = u'(x)\,\dd x$\,'' is nothing more
than a \emph{name} for this theorem. The $\dd x$ does not cancel
because of algebra; it cancels because the theorem says the two
integrals are equal.

Everything that looks like fraction juggling --- $u$-substitution,
separable ODEs, even the chain rule written as $\dfrac{\dd y}{\dd
x}=\dfrac{\dd y}{\dd u}\dfrac{\dd u}{\dd x}$ --- is this one fact
wearing different clothes. Leibniz deliberately designed the symbols
so that the theorem would \emph{look} like cancellation. That is a
feature, not an accident (more on this at the end).

A second, equivalent phrasing is worth keeping in your pocket,
because it makes the ODE argument airtight:

\begin{center}
  \itshape
  Two functions of $x$ with the same derivative differ only by a constant.
\end{center}

\section{Separable ODEs: the shorthand, then the expansion}

A separable equation has the form
\[
  \ddx{y} \;=\; g(x)\,h(y).
\]

\textbf{Shorthand (the way it is usually written).}
Treat $\dfrac{\dd y}{\dd x}$ as a fraction and shuffle symbols:
\[
  \frac{\dd y}{h(y)} \;=\; g(x)\,\dd x
  \qquad\Longrightarrow\qquad
  \int \frac{\dd y}{h(y)} \;=\; \int g(x)\,\dd x .
\]
Two genuinely forbidden moves were used: ``multiplying by $\dd x$''
and ``slapping $\int$ on a loose $\dd y$.''

\textbf{Expanded (what the shorthand abbreviates).}
No symbol is ever split. Remember that $y$ is an unknown
\emph{function} $y(x)$, and assume $h(y)\neq 0$ on the interval of interest.

\textbf{Step 1.} Divide both sides by $h(y(x))$ --- legal, ordinary
division of numbers:
\[
  \frac{1}{h(y(x))}\,\ddx{y} \;=\; g(x).
\]

\textbf{Step 2.} Both sides are now ordinary functions of $x$. If
they are equal, their antiderivatives in $x$ are equal up to a constant:
\[
  \int \frac{1}{h(y(x))}\,\ddx{y}\,\dd x \;=\; \int g(x)\,\dd x .
\]

\textbf{Step 3.} Apply the substitution rule to the left integral,
with the inner function $u=y(x)$ and $f(y)=1/h(y)$:
\[
  \int \underbrace{\frac{1}{h(y(x))}\,\ddx{y}\,\dd
  x}_{\text{theorem: this equals}\ \int \frac{1}{h(y)}\,\dd y}
  \;=\; \int \frac{1}{h(y)}\,\dd y .
\]

Combining the three steps gives \emph{exactly} the shorthand's conclusion,
\[
  \int \frac{1}{h(y)}\,\dd y \;=\; \int g(x)\,\dd x,
\]
but nothing was ever divided into ``$\dd y$ over $\dd x$.'' The
shorthand and the expansion produce identical work; the difference is
only that the expansion names a theorem at each step instead of
pretending the symbol is a fraction.

\section{A worked example, side by side}

Solve $\;\ddx{y} = x\,y.\;$ Here $g(x)=x$ and $h(y)=y$.

\textbf{Shorthand.}
\[
  \begin{aligned}
    \ddx{y} &= x\,y \\[3pt]
    \frac{1}{y}\,\dd y &= x\,\dd x \\[3pt]
    \int \frac{1}{y}\,\dd y &= \int x\,\dd x \\[3pt]
    \ln|y| &= \tfrac{x^2}{2} + C \\[3pt]
    y &= A\,e^{x^2/2}
  \end{aligned}
\]

\textbf{In full detail.}
$y=y(x)$, and suppose $y\neq 0$:
\[
  \begin{aligned}
    \frac{1}{y(x)}\,y'(x) &= x \\[3pt]
    \frac{\dd}{\dd x}\ln|y(x)| &= \frac{\dd}{\dd
    x}\!\left(\tfrac{x^2}{2}\right) \\[3pt]
    \therefore\quad \ln|y(x)| &= \tfrac{x^2}{2} + C \\[3pt]
    y &= A\,e^{x^2/2}
  \end{aligned}
\]

In the rigorous column, the second line is the \emph{whole trick}: by
the chain rule $\dfrac{\dd}{\dd x}\ln|y(x)| = \dfrac{y'(x)}{y(x)}$,
so the left side really is the derivative of $\ln|y|$. Once both
sides are recognised as derivatives of known functions, ``two
functions with the same derivative differ by a constant'' finishes
the job. The shorthand's ``$\int \tfrac1y\,\dd y = \ln|y|$'' was
secretly this chain-rule recognition all along.

\textit{An honest footnote the shorthand hides:} dividing by $y$ in
step one quietly assumed $y\neq 0$, so it \emph{drops} the constant
solution $y\equiv 0$. The expanded constant $A=\pm e^{C}$ ranges over
nonzero values; allowing $A=0$ afterwards restores the lost solution.
The careful version makes such assumptions visible.

\section{The same engine drives $u$-substitution}

Evaluate $\displaystyle\int 2x\cos(x^2)\,\dd x$. The shorthand sets
$u=x^2$, writes ``$\dd u = 2x\,\dd x$,'' and gets
\[
  \int \cos u \,\dd u = \sin u + C = \sin(x^2)+C.
\]
Why is ``$\dd u = 2x\,\dd x$'' allowed? Because it is the same
substitution theorem, and you can verify the result with no
differentials at all: by the chain rule,
\[
  \frac{\dd}{\dd x}\sin(x^2) = \cos(x^2)\cdot 2x,
\]
which is exactly the integrand. The ``$\dd u$'' bookkeeping is just a
fast way to organise this chain-rule check.

\section{Exact equations: the same trick in two variables}

Consider a curve defined implicitly by $F(x,y)=C$, describable locally as
$y=y(x)$. Differentiating both sides looks, on its face, like more
forbidden fraction algebra:
\[
  F_x + F_y\,\ddx{y} = 0
  \qquad\Longrightarrow\qquad
  F_x\,\dd x + F_y\,\dd y = 0,
\]
where $F_x=\pd{F}{x}$ and $F_y=\pd{F}{y}$ are evaluated along the curve.
The right-hand equation is exactly the form
\[
  M\,\dd x + N\,\dd y = 0, \qquad M=F_x,\ N=F_y,
\]
that a differential-equations course calls an \textbf{exact equation}.
The step from the implicit equation to this one again looks like
cross-multiplying $\dd x$ into the indivisible symbol $\dfrac{\dd y}{\dd
x}$. What actually licenses it?

\subsection{The honest derivation}

\begin{theorem}[Total derivative, i.e.\ the chain rule for two variables]
  Let $F(x,y)$ have continuous partial derivatives, and let $y=y(x)$
  be differentiable. Then
  \[
    \frac{\dd}{\dd x}F\bigl(x,y(x)\bigr) \;=\; F_x\bigl(x,y(x)\bigr) +
    F_y\bigl(x,y(x)\bigr)\,y'(x).
  \]
\end{theorem}
This is proved the same way the ordinary chain rule is --- by bounding
the error in the linear approximation of $F$ near a point and taking a
limit --- with no differential ever divided or multiplied on its
own\footnote{The full $\varepsilon$--$\delta$ argument is carried out
in the companion note \emph{Rigorous Derivation of the Total
Derivative}.}. It is an honest statement about the single real-valued
function $x\mapsto F(x,y(x))$.

Now suppose $F(x,y(x))=C$ for all $x$ in some interval --- that is
what ``$F(x,y)=C$ defines $y(x)$'' means. Differentiating the constant
function gives $0$, so by the theorem,
\[
  F_x + F_y\,y'(x) = 0
\]
for every such $x$. This is Step~1 of the separable-ODE argument again,
in the same disguise: one honest equation between functions of $x$,
obtained with the chain rule and never by splitting $\dd y/\dd x$.

\subsection{What ``$M\,\dd x+N\,\dd y=0$'' really packages}

Run the logic forward to see what the compact notation stands for.
Define, for a function of two variables, its \textbf{total
differential}
\[
  \dd F \;:=\; F_x\,\dd x + F_y\,\dd y .
\]
Just as $\dd y = \ddx{y}\,\dd x$ is a genuine identity between
$1$-forms (Section~\ref{sec:deeper}), $\dd F$ is not built by dividing
anything --- it is a new object, a linear combination of the symbols
$\dd x$ and $\dd y$, defined precisely so that its restriction to a
curve $y=y(x)$ recovers the chain-rule expression above:
\[
  \dd F \big|_{y=y(x)} \;=\; \bigl(F_x + F_y\,y'(x)\bigr)\,\dd x .
\]
So the sentence ``$F_x\,\dd x+F_y\,\dd y=0$'' is shorthand for ``the
total derivative of $F$ along the curve is zero,'' which by the
theorem above and \emph{equal derivatives differ by a constant}
(applied to the identically-zero derivative) means $F(x,y(x))$ is
constant --- exactly the $C$ we started with. Nothing was
cross-multiplied; the notation was engineered so that packaging the
chain rule as ``$\dd F=0$'' and unpacking it back to ``$F=C$'' are
inverse operations.

\subsection{Running it the other way: recovering $F$ from $M$ and $N$}

The differential-equations use of this idea starts from the other end.
One is handed
\[
  M(x,y)\,\dd x + N(x,y)\,\dd y = 0
\]
with no $F$ in sight, and asked to solve it. The equation deserves to
be called \emph{exact} precisely when some $F$ exists with $F_x=M$ and
$F_y=N$ --- i.e., when $M\,\dd x+N\,\dd y$ really is $\dd F$ for some
function $F$, so the argument above applies and the solution curves
are the level sets $F(x,y)=C$.

\begin{proposition}[Exactness test]
  If $M=F_x$ and $N=F_y$ for some twice continuously differentiable
  $F$, then
  \[
    \pd{M}{y} = F_{xy} = F_{yx} = \pd{N}{x},
  \]
  by equality of mixed partial derivatives. Conversely, if $M_y=N_x$
  throughout a simply connected region, such an $F$ exists.
\end{proposition}
The forward direction is Clairaut's theorem in disguise --- another
``obvious'' cancellation ($F_{xy}=F_{yx}$) that is really a
nontrivial theorem about mixed partials, not a symbol-swap. The
converse is the same existence fact behind Section~\ref{sec:deeper}'s
$1$-forms: a $1$-form with zero exterior derivative is locally exact
(the Poincar\'e lemma). The condition $M_y-N_x=0$ is precisely the
statement that $\dd(M\,\dd x+N\,\dd y)=0$.

\textbf{Worked example.} Solve $(2xy+3)\,\dd x+(x^2-1)\,\dd y=0$. Here
$M_y=2x=N_x$, so the equation is exact. Integrate $M$ in $x$, holding
$y$ fixed:
\[
  F(x,y) = \int (2xy+3)\,\dd x = x^2y+3x+g(y).
\]
Match $F_y$ against $N$ to pin down $g$:
\[
  F_y = x^2+g'(y) \overset{!}{=} x^2-1
  \quad\Longrightarrow\quad
  g'(y)=-1 \quad\Longrightarrow\quad g(y)=-y.
\]
So $F(x,y)=x^2y+3x-y$, and the implicit solution is
\[
  x^2y+3x-y=C.
\]
Differentiating this back with the ordinary and chain rules reproduces
the original equation exactly --- confirming that ``cross-multiplying
by $\dd x$'' and ``taking $\dd$ of both sides of $F=C$'' are the same
theorem viewed from opposite ends.

\medskip
\textit{A caution.} Do not confuse this additive total differential
with the multiplicative cyclic relation of
Section~\ref{sec:fails}. $\dd F=F_x\,\dd x+F_y\,\dd y$ is a
\emph{sum}, licensed by the chain rule for a function of two
variables; $(\partial x/\partial y)(\partial y/\partial
z)(\partial z/\partial x)=-1$ is a \emph{product} of three different
partials of three different two-variable functions implicitly linked
by one relation. No amount of total-differential bookkeeping turns one
into the other. Keeping straight \emph{which} theorem is in play is,
as always, the whole game.

\section{Where the fraction picture genuinely fails}
\label{sec:fails}

The cancellation only works for a \emph{first} derivative of
\emph{one} variable. Push it further and it bites back:

\textbf{Second derivatives don't cancel.} The symbol $\dfrac{\dd^2
y}{\dd x^2}$ is not $(\dd y)^2/(\dd x)^2$; you cannot treat the
``$2$''s as exponents and simplify. It is shorthand for
$\dfrac{\dd}{\dd x}\!\left(\dfrac{\dd y}{\dd x}\right)$.

\medskip
\textbf{Partial derivatives don't cancel.} For three quantities
linked by one relation,
\[
  \left(\frac{\partial x}{\partial y}\right)\!\left(\frac{\partial
  y}{\partial z}\right)\!\left(\frac{\partial z}{\partial x}\right) = -1,
\]
\emph{not} $+1$. Naive ``cancellation'' gives the wrong sign --- a
famous trap in thermodynamics.

This is exactly why instructors warn you off the fraction picture: it
is a reliable mnemonic in a narrow setting and a source of real
errors outside it. Knowing \emph{which theorem} you are using tells
you precisely when the shortcut is safe.

\section{The deeper resolution: differentials really are objects}
\label{sec:deeper}

There is a more advanced viewpoint in which $\dd x$ and $\dd y$ are
not informal shorthand at all, but well-defined mathematical objects
called \textbf{differential $1$-forms}. In that framework, for a
function $y$ of a single variable $x$, the equation
\[
  \dd y \;=\; \ddx{y}\,\dd x
\]
is a genuine, provable identity --- a true statement about $1$-forms,
not a mnemonic. Integration $\int(\cdots)\,\dd x$ becomes integration
of forms, and the substitution rule becomes the natural ``change of
variables'' for forms.

So Leibniz's notation was not merely lucky. It anticipated, by about
two centuries, a rigorous theory in which the symbols mean something
and the cancellation is a theorem. Introductory courses say ``not a
fraction'' because that theory is out of reach at that stage --- but
the manipulations were never wrong, only unjustified \emph{at that level}.

\textbf{In one breath.}
$\dfrac{\dd y}{\dd x}$ is a limit, not a fraction --- so the
warning is true. But every fraction-style manipulation in
$u$-substitution and separable ODEs is a disguised application of
the substitution rule (the chain rule backwards), together with
``equal derivatives differ by a constant.'' Exact equations run the
same trick with the multivariable chain rule: $F_x\,\dd x+F_y\,\dd
y=0$ is shorthand for $\dd F=0$, and $M_y=N_x$ is shorthand for
``a potential function $F$ exists,'' which is Clairaut's theorem
and the Poincar\'e lemma in disguise. Leibniz's notation is
engineered so those theorems \emph{look} like cancellation, and at
the level of differential $1$-forms the cancellation becomes
literally true. Trust the shortcut for first derivatives in one
variable, and for total differentials built additively from partial
derivatives; reach for the theorem the moment you hit higher
derivatives or naively cancelled partials.

\end{document}
